\documentclass{article}
\usepackage{/home/user1/code/tex/sty}
\usepackage{amsfonts}
\usepackage{amsmath}
\usepackage{cancel}
\usepackage{geometry}
\setlength{\parindent}{0pt}
\usepackage[backend=bibtex]{biblatex}
\addbibresource{bib.bib}
\geometry{top=1in,bottom=1in,left=1in,right=1in}
\begin{document}
\begin{center}
{\Large{Problem Set 3}}\\~\\
\begin{tabular}{l l}
Student&Agustin Romero\\
Instructor&Professor Savas Dimopoulos\\
Assistants&Jed Thomson\\
Course&Standard Model\\
Institution&Stanford University\\
Date&\today\\
\end{tabular}
\end{center}
\section*{Problem 1.a.}
Start with the $SU(2)\times U(1)$ interaction, insert a potential, then derive the $\cos(\theta_w)$ term.
\e{\phi=2^{-1/2} \begin{pmatrix} v \\ v \end{pmatrix}}
\e{\label{aodsf9}(D_\mu \phi)^\dagger (D^\mu \phi)&= |(-ig \oot \vec{\tau} \vec{W}_\mu -ig' \oot B_\mu)\phi |^2\\
&=\fr{1}{8} |\begin{pmatrix} gW_\mu^3+g' B_\mu & g(W_\mu^1 -iW_\mu^2) \\ g(W_\mu^1 + iW_\mu^2) & -gW_\mu^3+g'B_\mu \end{pmatrix} \begin{pmatrix} v \\ v \end{pmatrix}|^2\\
&= \fr{1}{8}(gv)^2((W_\mu^1)^2+(W_\mu^2)^2)+\fr{1}{8}v^2\begin{pmatrix} W_\mu^3 & B_\mu \end{pmatrix}\begin{pmatrix} g^2 & -gg' \\ gg' & g'^2 \end{pmatrix} \begin{pmatrix} W^{3\mu} \\ B^\mu \end{pmatrix}}
\e{m_W=\oot v g}
\e{m_z=\oot v(g+g'^2)^{1/2}}
\e{\boxed{\fr{m_W}{m_Z}=g(g^2+g'^2)^{-1/2}=\cos(\theta_w)}}
By expanding electroweak interaction terms in the Lagrangian we can derive the tree-level masses of $m_W$ and $m_Z$, and $\cos(\theta_w)$.
\section*{Problem 1.b.}
Add the new operator and expand.
\e{\fr{1}{\Lambda_T^2}(h^\dagger D_\mu h)(D^\mu h^\dagger h)&=\fr{1}{8 \Lambda^2_T}h^\dagger|\begin{pmatrix} gW_\mu^3+g' B_\mu & g(W_\mu^1 -iW_\mu^2) \\ g(W_\mu^1 + iW_\mu^2) & -gW_\mu^3+g'B_\mu \end{pmatrix} \begin{pmatrix} v \\ v \end{pmatrix}|^2 h\\
&= \boxed{\fr{g^2 v^4}{8 \Lambda_T^2}((W_\mu^1)^2+(W_\mu^2)^2)+\fr{v^4}{8 \Lambda_T^2}v^2\begin{pmatrix} W_\mu^3 & B_\mu \end{pmatrix}\begin{pmatrix} g^2 & -gg' \\ gg' & g'^2 \end{pmatrix} \begin{pmatrix} W^{3\mu} \\ B^\mu \end{pmatrix}}}
By expanding the new interaction term at the minimum of the Higgs potential the term is essentially the $O(v^4)$ electroweak interaction terms for W and Z with some coupling strength $1/\Lambda^2_T$.
\section*{Problem 1.c.}
Start with the Z current $j_Z$ in terms of hypercharge current $j_Y$ and weak neutral current $j_3$, then manipulate the current until it is in terms of $j_3$ and $j_{em}$. We should expect $j_3$ to couple only left handed particles, and $j_{em}$ ot couple left and right handed particles identically.
\e{\label{81n2}g_Z=\fr{g_W}{\cos(\theta_W)}}
\e{j_Z^\mu&=-j^\mu_Y\sin(\theta_w)+j_3^\mu \cos(\theta_w)\\
&=-\oot g' \sin(\theta_w)(Y_{f_L}\bar{u}_L \gamma^\mu u_L + Y_{f_R} \bar{u}_R \gamma^\mu u_R)+I_W^{(3)} g_W \cos(\theta_w) \bar{u}_L \gamma^\mu u_L\\
&=-\oot g' \sin(\theta_w)((Q_f- I_W^{(3)})\bar{u}_L \gamma^\mu u_L + Q_f \bar{u}_R \gamma^\mu u_R)+I_W^{(3)} g_W \cos(\theta_w) \bar{u}_L \gamma^\mu u_L\\
&=(-g' (Q_f-I_W^{(3)}) \sin(\theta_w)+I_W^{(3)} g_W \cos(\theta_w))\bar{u}_L \gamma^\mu u_L - g' \sin(\theta_w) Q_f \bar{u}_R \gamma^\mu u_R\\
&=g_w(-(Q_f-I_W^{(3)}) \fr{\sin^2(\theta_w)}{\cos(\theta_w)}+I_W^{(3)} \cos(\theta_W)) \bar{u}_L \gamma^\mu u_L - g_W \fr{\sin(\theta_w)}{\cos(\theta_W)} Q_f \bar{u}_R \gamma^\mu u_R
\intertext{Use \eref{81n2}.}
&=g_Z(I_W^{(3)}-Q_f \sin^2(\theta_w)) \bar{u}_L \gamma^\mu u_L -g_ZQ_f \sin^2(\theta_W) \bar{u}_R \gamma^\mu u_R
\intertext{Rearrange.}
&=g_Z I_W^{(3)} \bar{u}_L \gamma^\mu u_L - g_ZQ_f \sin^2(\theta_W) (\bar{u}_L \gamma^\mu u_L \bar{u}_R \gamma^\mu u_R)
\intertext{The first term is the $SU(2)_L$ current, and the is the electromagnetic current.}
&=\boxed{\fr{g_W}{\cos(\theta_W)} j_L^\mu - \sin^2(\theta_W) j_{em}^\mu}}
\section*{Problem 2.a.}
Start with the QED Lagrangian, compute the conserved current $j$, then the conserved charge $Q$, then show the charge is invariant under interactions. The Lagrangian for QED is
\e{D_\mu=\p_\mu+ieA_\mu(x)}
\e{\scl&=\bar{\psi} (i\cancel{\p}-m)\psi-\oof F_{\mu\nu}F^{\mu\nu}-e\bar{\psi}\gamma^\mu \psi A_\mu\\
&=\bar{\psi}(i\cancel{D}-m)\psi-\oof F_{\mu\nu} F^{\mu\nu}}
Its Noether current is
\e{ej^\nu=e\bar{\psi}\gamma^\nu \psi=\p_\mu F^{\mu\nu}}
Its Noether charge is
\e{Q=\int j^0 d^3x = \int \bar{\psi}(x) \psi(x) d^3x}
Under a local phase transofmration $\theta(x)$, $Q$ is invariant.
\e{\boxed{Q(\psi)\rightarrow Q(e^{i\theta(x) }\psi)=\int \bar{\psi}(x) e^{i\theta(x)} e^{-i\theta(x)} \psi(x) d^3x = Q}}
\section*{Problem 2.b.a.}
The QED $U(1)$ phase transformation is
\e{\boxed{\psi(x) \rightarrow e^{-i\theta(x)}\psi}}
The conserved $U(1)$ current is
\e{j^\mu(x)=\bar{\psi}(x)\gamma^\mu \psi(x)}
The symmetry is \fbox{not broken}.
\section*{Problem 2.b.b.}
Refer to Donoghue p77. The $SU(2)_\rom{isospin}$ group transforms the $u$ and $d$ fields in this way,
\e{\boxed{\begin{pmatrix}u'\\d'\end{pmatrix}=U\begin{pmatrix}u\\d\end{pmatrix}=e^{i\vec{\theta} \vec{\tau}} \begin{pmatrix}u\\d\end{pmatrix}}}
where $\vec{\theta}$ is the transformation parameter and $\vec{\tau}=\fr{1}{2}\vec{\sigma}$. There are 3 Noether currents, 
\e{\boxed{j_\mu^i=\bar{\psi}\gamma_\mu \oot \sigma^i \psi}}
where $\mu$ is a Dirac spinor index, and $i$ is an isospin index.
The $u$ and $d$ fields are associated with different electromagnetic charges and masses. $SU(2)_\rom{isospin}$ transformations do not conserve these charges or masses, so the \fbox{$SU(2)_\rom{isospin}$ is not conserved nor spontaneously broken}. There are combinations of fields that do conserve isospin charge, for example
\e{2^{-1/2}(ud-du)}
which is a singlet under $SU(2)$ with isospin charge 0.
\section*{Problem 2.b.c.}
Refer to Donoghue p103. The transformation for $B$ charge is
\e{\boxed{q\rightarrow e^{i\phi_B} q}}
The Noether current is
\e{\boxed{j^\mu = \oot (\bar{u}\gamma^\mu u+\bar{d}\gamma^\mu d+...)}}
B charge is \fbox{not conserved} due to triangle diagrams involving vector-vector-axial vertices.
\section*{Problem 2.b.d.}
Refer to Donoghue p103. The transformation for $L$ charge is
\e{\boxed{l\rightarrow e^{i\phi_L} l}}
The Noether current is
\e{\boxed{j^\mu = \bar{e}\gamma^\mu e + \bar{\nu}_{eL} \gamma^\mu \nu_{eL} + ...}}
L charge is \fbox{not conserved} due to triangle diagrams involving vector-vector-axial vertices.
\section*{Problem 2.b.e.}
$B-L$ is the combination of the above $B$ and $L$ transformations, thus applicable to lepton-quark interactions. The currents show that
\e{\p_\mu (j^\mu_B-j^\mu_L)=0}
Therefore $B-L$ charge is \fbox{conserved}.
\section*{Problem 3.a.}
Starting with the Higgs kinetic term,
\e{\fr{1}{8} |\begin{pmatrix} gW_\mu^3+g' B_\mu & g(W_\mu^1 -iW_\mu^2) \\ g(W_\mu^1 + iW_\mu^2) & -gW_\mu^3+g'B_\mu \end{pmatrix} \phi|^2}
the number of Higgs doublets in $\phi$ can be arbitrary, and so
\e{\boxed{m_Z^2=\fr{m_W^2}{\cos(\theta_w)}}}
holds for any number of Higgs doublets \cite{LANGACKER1981185}.
\section*{Problem 3.b.}
Let $n$ be the number of scalar fields, $N$ the number of generators of a Lie Group $G$, $M$ the number of generators that leave the vacuum $|0\rangle$ invariant, and $N-M$ the number of generators that do not. Then there are $p=N-M$ number of spontaneously broken symmetries, and $p$ number of massive Goldstone Bosons. The arbitrary number of Higgs fields $\phi$ is
\e{\phi\rightarrow U\phi = v+\eta+=\begin{pmatrix} v_1 + \eta_1 \\ \vdots \\ v_p + \eta_p \\ 0 \\ \vdots \\ 0\end{pmatrix}}
Then expand the kinetic term for these Higgs fields, similar to \eref{aodsf9},
\e{\oot D^\mu \phi D_\mu \phi &= (v+\eta)^T(p^\mu + ig\vec{A}^\mu\cdot \vec{L})(\p_\mu - ig \vec{A}_\mu \cdot \vec{L})(v+\eta)\\
&=\oot v^T L^i L^j v A^{i\mu} A^j_\mu + \rom{interaction terms}\\
&=\oot M_{ij}^2 A^{i\mu} A^j_\mu}
where there is now the mass matrix $M_{ij}$,
\e{M_{ij}^2=M_{ji}^2=g^2 v^T L^iL^j v = g^2 \langle v | L^i L^j v \rangle = g^2 \langle L^i v | L^j v \rangle}
The masses of the W and Zs are therefore
\e{M_Z^2=(g^2+g'^2)\Sigma_N (t_n^3)^2 v_n^2}
\e{M_W^2=\oot g^2 (t_n(t_n+1)-(t_n^3)^2)v_n^2}
Take the ratio of $M_W$ and $M_Z \cos(\theta_w)$,
\e{\label{9128n}\rho=\fr{M_W^2}{\cos(\theta_W)M_Z^2}=\fr{\Sigma (t_n(t_n+1)-(t_n^3)^2)v_n^2}{2\Sigma_n (t_n^3)^2 v_n^2}}
where $T$ is the isospin charge, $Y$ is the hypercharge, and $v$ is the vacuum expectation value of the $i$-th Higgs field.
\section*{Problem 3.c.}
The constraint for $\rho=1$ on \eref{9128n} implies
\e{4T(T+1)-Y^2=2Y^2}
Simplifying,
\e{\boxed{T(T+1)=\fr{3}{4}Y^3}}
There is a hyperbola of solutions, where
\e{\boxed{Y=\pm \fr{2}{3^{1/2}} (T^2+T)^{1/2}}}
Which is satisfied when $T=\oot$, $Y=1$ as stated in the prompt. Another non-trivial solution is \fbox{$T=3$, $Y=4$}.
\printbibliography
\end{document}
